%==============================================================================
% Section 5: Discussion
% Verbatim from Nckh-official_2.docx (P86-P96).
%==============================================================================
\section{Discussion}
\label{sec:discussion}

The 1.52\% MAE of the long-sequence model on the strict cross-battery protocol
indicates that the selective state space mechanism is able to capture long-range
degradation dependencies without the recurrence bottleneck of LSTM-based models
or the quadratic cost of Transformers~\cite{gu2023,meng2026,zhao2024}.

Our observation that attention pooling surpasses last-token aggregation aligns
with the ablation analysis in Section~\ref{subsec:ablation}, indicating that the
signatures of degradation are spread out instead of being localised in a single
terminal state during the discharge process, consistent with recent
sequence-learning studies~\cite{meng2026,nie2023} and confirming our design
choice for combining patch encoding with global token aggregation to maintain
local and long-range degradation properties.

% Table 6 of the source manuscript.
\begin{table}[H]
  \caption{Comparison with Recent Cross-Battery Studies on the NASA Dataset}
  \label{tab:comparison}
  \centering
  \scriptsize
  \linespread{0.92}\selectfont
  \tabtight
  % Fixed column widths sized to the widest unbreakable words ("Architecture",
  % "SambaMixer-", "Parameters"); the padding is trimmed so the X columns keep
  % enough measure.
  \setlength{\tabcolsep}{2pt}%
  \begin{tabularx}{\textwidth}{|M{1.7cm}|M{1.8cm}|M{1.6cm}|Y|Y|Y|M{1.6cm}|}
    \hline
    \textbf{Work} & \textbf{Architecture} & \textbf{Parameters} &
    \textbf{Protocol} & \textbf{MAE} & \textbf{Latency} &
    \textbf{Robustness} \\
    \hline
    SambaMixer-L (2025)
      & MambaMixer & 48.7M & 10 train / 3 eval & 0.51--1.20\%
      & Not reported & None \\
    \hline
    TIDSIT (2025)
      & Transformer & Not reported & 2 train / 1 eval & 0.58\%
      & Not reported & None \\
    \hline
    Long model (This work)
      & Mamba + FiLM & 99k & 23 train / 2 val / 1 test & 1.52\%
      & Not measured (see window-30) & LOBO (16 folds) \\
    \hline
    Window-30 model (This work)
      & Mamba + FiLM & 79k & Same as above & 1.98\%
      & 56.1\,ms (P95: 78.3\,ms) & --- \\
    \hline
  \end{tabularx}
  \srcnote{Source: Authors' processing, 2026}
\end{table}

The proposed framework does not depend on the injection of electrochemical
knowledge using electrochemical impedance spectroscopy (EIS) as was done in
previous multimodal Mamba studies~\cite{meng2026} but instead extracts FFT-based
spectral descriptors directly from discharge telemetry via FiLM conditioning,
thus bypassing the need for the specialised measurement hardware required by
EIS, while still taking advantage of the higher degradation sensitivity of
frequency-domain indicators compared to raw time-domain
signals~\cite{alsmadi2025,yang2023}.

The data-centric improvement shown in Section~\ref{subsec:main-results}, where
the addition of a single additional \degC{4} battery reduced MAE without any
architectural modification, reinforces the notion that operating-condition
coverage might be more important than model complexity for cross-battery
generalisation, consistent with prior research indicating that performance
degradation in battery prognostics is often caused by distribution shift rather
than insufficient network capacity~\cite{meng2026,yang2023}.

This accuracy--latency trade-off compact configuration for real-time edge
alerts, long-sequence configuration for periodic cloud diagnostics
(Section~\ref{subsec:latency}) is analogous to recent comparisons between Mamba
and Transformer architectures for resource-constrained
deployment~\cite{olalde2024}.

As per Section~\ref{subsec:anomaly}, the tiny F1 obtained using proxy labels is
enough for Isolation Forest as a complementary early-warning module rather than
an autonomous fault classifier, a limitation shared by unsupervised anomaly
detection in general where proxy labels cannot fully represent real operational
failures~\cite{aljohani2026,arbaoui2024}.

Three limitations need to be recognised: (i) validation was limited to the NASA
dataset (limited cells, single chemistry); (ii) despite the more realistic
evaluation, large LOBO performance variation persisted for batteries with rare
degradation trajectories; (iii) anomaly detection was based on proxy labels
without access to actual fault labels. Therefore, future work should include
validation on larger multi-chemistry datasets and real BMS fault logs, as well
as exploring semi-supervised or continual learning approaches for better
robustness.
